\documentclass[12pt]{article}
\usepackage[utf8]{inputenc}
\usepackage[T1]{fontenc}
\usepackage{hyperref}
\title{Yatay Language Resources for K-12 Curriculum Development}
\author{Compiled by Claude Sonnet 5.5 (claude-sonnet-5-5) for Andr\'as Kornai}
\date{}
\begin{document}
\maketitle

\section{Introduction}
Yatay (ISO 639-3: yty) is described in the reference material as an extinct Australian language. The search found no dedicated page, grammar, word list or teaching material for it, and could not independently confirm the code. The only related result is an Endangered Languages Project entry that appears to list Yatay among the ``also known as'' names of Yadhaykenu, a language of Cape York. That entry gives it as critically endangered, with very few speakers. It is not verified that this is the same language as yty, so treat the link as tentative.

\section{Primary Educational Resources}
None found.

\section{Linguistic and Reference Materials}
Endangered Languages Project entry (``Also Known As'' list that includes Yatay; the identification with yty is unconfirmed): \url{https://endangeredlanguages.com/node/38890}

\section{Religious and Community Resources}
None found.

\section{Digital and Audio Resources}
None found.

\section{Curriculum Development Context}
There is nothing in the search results to base a K-12 curriculum on. A next step would be to check AIATSIS holdings and the Austlang database for the correct language entry, then contact the relevant Aboriginal language centre. Resources should be added only after the identity of the language is confirmed.
\end{document}